\section{Conclusion}

Hardware design gives an agent benchmark a strong verifier and no statement of intent. Everything in copperbench follows from taking both halves of that seriously: verification is used where it is authoritative, end-state assertions supply what it cannot say, and the two are graded separately so that a legal design and a correct one are never confused.

Three choices are worth carrying into other agentic domains with the same shape. Grading offline and model-free makes a published number checkable by parties who do not trust the authors, and costs little once the runs write structured artifacts. A closed grading vocabulary trades expressive power for comparability and reviewability, and the trade is favorable as soon as a suite has more than a few authors. Scoring refusal as an outcome class, with a citation requirement and a task mix in which neither compliance nor caution wins, turns a safety property into something measurable rather than asserted.

The benchmark's most useful output is not the leaderboard. It is the ranked failure taxonomy: an ordered list of what to fix, priced by what fixing it is worth. We expect the scores here to age quickly. We expect the standard, and the discipline of generating every reported number from released evidence, to age better.
